\documentclass[11pt]{article}
\usepackage[review]{acl}
\usepackage{times}
\usepackage{latexsym}
\usepackage[T1]{fontenc}
\usepackage[utf8]{inputenc}
\usepackage{microtype}
\usepackage{inconsolata}
\usepackage{booktabs}
\usepackage{amsmath}
\usepackage{amssymb}

% Generated from results/paper_stats.json; do not edit.
\newcommand{\ModelName}{NLLB-200 distilled 600M}
\newcommand{\SyntheticStringCount}{3}
\newcommand{\SyntheticOutcomeCount}{5}
\newcommand{\SyntheticBatchSize}{2}
\newcommand{\SyntheticMCBatches}{200,000}
\newcommand{\QLow}{0.25}
\newcommand{\QHigh}{0.75}
\newcommand{\ResponsePhiLow}{0.0}
\newcommand{\ResponsePhiHigh}{0.0}
\newcommand{\GlobalPhiLow}{0.01198}
\newcommand{\GlobalPhiHigh}{0.01906}
\newcommand{\ResponseThetaMaxDiff}{5.55e-17}
\newcommand{\MCResponseThetaError}{2.04e-04}
\newcommand{\MCResponsePhiError}{1.56e-04}
\newcommand{\MCGlobalThetaError}{4.28e-05}
\newcommand{\MCGlobalPhiError}{2.13e-06}
\newcommand{\TatoebaSources}{975}
\newcommand{\TatoebaOutputs}{1,950}
\newcommand{\CalibrationReplicates}{2}
\newcommand{\TiltBeta}{0.20}
\newcommand{\TiltTemperature}{0.70}
\newcommand{\MatchedTemperature}{0.70}
\newcommand{\EntropyTolerance}{0.05}
\newcommand{\EntropyResidual}{-0.0099}
\newcommand{\FertilityDelta}{+0.031}
\newcommand{\FertilityLow}{+0.020}
\newcommand{\FertilityHigh}{+0.043}
\newcommand{\ChrfDelta}{-0.021}
\newcommand{\ChrfLow}{-0.032}
\newcommand{\ChrfHigh}{-0.010}
\newcommand{\BleuDelta}{-2.20}
\newcommand{\BleuLow}{-3.84}
\newcommand{\BleuHigh}{-0.62}
\newcommand{\ActionDelta}{+0.146}
\newcommand{\ActionLow}{+0.068}
\newcommand{\ActionHigh}{+0.225}
\newcommand{\WordDelta}{+0.010}
\newcommand{\WordLow}{-0.043}
\newcommand{\WordHigh}{+0.063}
\newcommand{\EOSCount}{1,950}
\newcommand{\TruncationCount}{0}
\newcommand{\EmptyCount}{0}
\newcommand{\SpecialTokenCount}{0}


\title{Exact Segmentation Invariance and a Negative Feasibility Result for Tokenization-Oriented Machine Translation Policy Gradients}

\author{Anonymous ACL submission}

\begin{document}
\maketitle

\begin{abstract}
We investigate whether changing target-side segmentation can provide a clean
intervention for reinforcement learning in machine translation.  The project
produced two complementary negative results.  First, for a policy that samples
a surface string and then samples a conditionally normalized segmentation,
ordinary response-level policy gradients with string-only reward are exactly
invariant to the segmentation sampler.  Token-normalized or position-weighted
updates can break this invariance, but then the effect belongs to the estimator.
Second, a normalized piece-length tilt in \ModelName{} changed the complete
autoregressive policy rather than merely resegmenting fixed strings.  On a
frozen English--German calibration set it increased token fertility while
reducing both the fixed string reward and an independent translation metric.
The prespecified launch screen was therefore not passed, and no research
training run was launched.  Together, these results clarify why exact behavior
likelihood is necessary but insufficient for identifying a useful
tokenization-only intervention in sequence-level reinforcement learning.
\end{abstract}

\section{Motivation}

Subword regularization and BPE-dropout show that sampling alternative
segmentations can improve supervised machine translation
\citep{kudo-2018-subword,provilkov-etal-2020-bpe}.  This does not immediately
imply a useful reinforcement-learning intervention.  Policy-gradient training
requires the likelihood of the process that generated an action, and changing
the segmentation may either leave a string-level estimator unchanged or alter
the represented autoregressive policy itself.  Exploration and reward scaling
are already known to matter in neural machine translation policy gradients
\citep{kiegeland-kreutzer-2021-revisiting}.

We report the completed tokenization branch of a larger project.  Our scope is
deliberately narrow: estimator validity and intervention feasibility.  We do
not report an effect of tokenization on reinforcement-learning outcomes because
the intervention failed its calibration gate before training.

\section{Same-string intervention}

Let a string policy sample a surface string $s$, after which a segmentation
sampler draws a path $z$:
\begin{equation}
P_{\theta,\phi}(s,z\mid x)
=P_\theta(s\mid x)q_\phi(z\mid s,x).
\end{equation}
For string-only reward $R(s)$ and a normalized conditional sampler, define the
two expected score-function gradients as
\begin{align}
g_\theta
&=\mathbb{E}_{s,z}[R(s)\nabla_\theta\log P_{\theta,\phi}(s,z\mid x)]\\
&=\mathbb{E}_{s}[R(s)\nabla_\theta\log P_\theta(s\mid x)],\\
g_\phi
&=\mathbb{E}_{z\mid s}[\nabla_\phi\log q_\phi(z\mid s,x)]
=\mathbf{0}.
\end{align}
Thus a separately sampled segmentation cannot change the expected ordinary
string-policy gradient.  Scoring the alternate token path with the base
model's raw autoregressive likelihood would instead score a process that did
not generate the sample.

\subsection{Exact finite verification}

We enumerate a categorical policy over \SyntheticStringCount{} strings and
\SyntheticOutcomeCount{} joint string--path outcomes.  Some strings admit
short and long paths; reward depends only on the string.  We enumerate ordered
batches of size \SyntheticBatchSize{} and compare exact expectations with
\SyntheticMCBatches{} Monte Carlo batches.

Table~\ref{tab:synthetic} varies the probability of the short path.  The
response-mean string gradient changes by at most \ResponseThetaMaxDiff{}, and
the segmentation-parameter gradient remains zero.  Dividing by the random
total number of tokens in each batch instead changes expected sample weighting
and yields a nonzero segmentation gradient.

\begin{table}[t]
\centering
\small
\begin{tabular}{lrr}
\toprule
Estimator & $q=\QLow$ & $q=\QHigh$ \\
\midrule
Response mean, $\phi$ gradient & \ResponsePhiLow & \ResponsePhiHigh \\
Global token mean, $\phi$ gradient & \GlobalPhiLow & \GlobalPhiHigh \\
\bottomrule
\end{tabular}
\caption{Exact segmentation-sampler gradients as the short-path probability
$q$ varies.  The global token mean has a sampled, path-dependent denominator.}
\label{tab:synthetic}
\end{table}

The maximum Monte Carlo error for the response-mean string-gradient
coordinates is \MCResponseThetaError{}; its segmentation-gradient error is
\MCResponsePhiError{}.  Corresponding errors for the global-token estimator
are \MCGlobalThetaError{} and \MCGlobalPhiError{}.  These checks validate the
implementation rather than provide translation-quality evidence.

\section{Autoregressive length tilt}

Because the same-string construction leaves no ordinary-gradient mechanism,
we also tested a broader decoder-policy intervention in \ModelName{}, a
released NLLB model \citep{nllb-2024-scaling}.  At each decoding step, a fixed
bias favored shorter visible SentencePiece tokens and the categorical policy
was renormalized.  The sampled token then changed the next decoder input,
position, attention state, future distribution, and stopping time.  This is a
total autoregressive policy intervention, not reassignment of a segmentation
to an otherwise fixed translation.

The frozen comparison used a tilt coefficient of \TiltBeta{} and temperature
\TiltTemperature{}.  The control used the entropy-matched temperature
\MatchedTemperature{}.  Entropy was measured on shared teacher-forced prefixes;
the residual was \EntropyResidual{} nats against a tolerance of
\EntropyTolerance{} nats.  The calibration contained \TatoebaSources{} unique
Tatoeba sources and \CalibrationReplicates{} paired rollout replicates, giving
\TatoebaOutputs{} generated outputs.  Human scores were not used.

\begin{table}[t]
\centering
\small
\begin{tabular}{lrrr}
\toprule
Tilt minus control & Mean & Lower & Upper \\
\midrule
Tokens per word & \FertilityDelta & \FertilityLow & \FertilityHigh \\
chrF++ reward & \ChrfDelta & \ChrfLow & \ChrfHigh \\
Corpus BLEU points & \BleuDelta & \BleuLow & \BleuHigh \\
Actions incl. EOS & \ActionDelta & \ActionLow & \ActionHigh \\
Decoded words & \WordDelta & \WordLow & \WordHigh \\
\bottomrule
\end{tabular}
\caption{Paired source-bootstrap intervals on the frozen calibration set.
The chrF++ reward is on its normalized scale; BLEU is reported in points.}
\label{tab:calibration}
\end{table}

Table~\ref{tab:calibration} shows the central feasibility result.  Fertility
increased, but the reward and independent metric both decreased.  The action
count changed while decoded word count was nearly stable, supporting a
tokenization shift, but not a competence-preserving one.  All \EOSCount{}
outputs ended in EOS; truncation, empty-output, and non-EOS special-token counts
were respectively \TruncationCount{}, \EmptyCount{}, and \SpecialTokenCount{}.
The prespecified quality screen remained inconclusive about whether both losses
were within acceptable limits after its sole bounded extension.  The pilot was
not launched.

\section{Discussion}

The two studies rule out different shortcuts.  Exact joint likelihood does not
make a same-string segmentation sampler causally active under an ordinary
string-reward estimator.  Conversely, passing alternate token identities
through an autoregressive decoder supplies a mechanism, but it also changes
the policy's initial competence and surface-string distribution.  A fertility
change under token-normalized credit assignment is an estimator effect; a
fertility change under the length tilt is a total decoder-policy effect.

This distinction also bounds the negative result.  We do not conclude that
tokenization is irrelevant to reinforcement learning.  We conclude that the
tested clean construction is invariant, while the tested non-invariant
construction did not satisfy its launch gate.

\section{Conclusion}

The tokenization branch is complete as a validation and feasibility study.
It provides an exact invariance result, tests where alternative reductions
break that invariance, and records a negative real-model calibration.  It does
not establish a translation-learning effect, because training was correctly
stopped before that estimand was evaluated.

\section*{Limitations}

The exact harness is intentionally finite and does not model neural decoder
states.  The real-model screen covers one model family, one language direction,
one calibration split, and one fixed tilt dose.  Entropy matching does not make
two policies equivalent.  The calibration estimates initial behavior only and
cannot determine whether later reinforcement learning would recover from the
quality loss.

\section*{Ethical Considerations}

The real-model study used the documented noncommercial research terms for the
NLLB checkpoint and attributed Tatoeba data.  No human judgments were newly
collected.  No model weights were released or trained after the feasibility
screen failed.

\bibliography{references}

\end{document}
